\documentclass[11pt]{article}
\usepackage[a4paper,margin=21mm]{geometry}
\usepackage{fontspec}
\setmainfont{Latin Modern Roman}
\newfontfamily\CJKfont{Noto Serif CJK SC}
\XeTeXlinebreaklocale "zh"
\XeTeXlinebreakskip=0pt plus 1pt
\usepackage{amsmath,amssymb,amsthm,amscd,graphicx,color}
\usepackage{xurl}
\usepackage[unicode,hidelinks]{hyperref}
\allowdisplaybreaks
\setlength\emergencystretch{3em}
\numberwithin{equation}{section}

\newtheorem{thm}{Theorem}[section]
\newtheorem*{Theorem*}{Theorem}
\newtheorem{cor}[thm]{Corollary}
\newtheorem{lmm}[thm]{Lemma}
\newtheorem{prp}[thm]{Proposition}
\newtheorem{mresult}[thm]{Main Result}
 {\theoremstyle{definition}
\newtheorem{que}[thm]{Question}
\newtheorem{conj}[thm]{Conjecture}
\newtheorem{eg}[thm]{Example}
\newtheorem{abuse_of_notation}[thm]{Abuse of Notation}
\newtheorem{claim}[thm]{Claim}
\newtheorem{verify}[thm]{Verification}
\newtheorem{result}[thm]{Result}

\newtheorem{notn}[thm]{Notation}
\newtheorem{exrc}[thm]{Exercise}

\newtheorem{rem}[thm]{Remark}
\newtheorem*{remm}{Remark}
\newtheorem{defn}[thm]{Definition} }
\makeatletter
\newcommand*{\rom}[1]{\expandafter\@slowromancap\romannumeral #1@}
\makeatother
\def \ran{\rangle}

\makeatletter\newlabel{m_rslt1}{{2.1}{}{Original source locator}{}{}}\makeatother
\makeatletter\newlabel{A1FT0T0_coll_cycle_ver}{{7.6}{}{Original source locator}{}{}}\makeatother
\begin{document}
\begin{center}\Large Tangency intersection cycles and marked-point coalescence\end{center}
\noindent\textbf{Stable ID:} 552\quad\textbf{Author:} Indranil Biswas; Apratim Choudhury; Ritwik Mukherjee; Anantadulal Paul\par
\noindent\textbf{Work:} Counting Curves with Tangencies\par
\noindent\textbf{Source:} \url{https://sigma-journal.com/2025/012/}\par
\noindent\textbf{Version:} Official SIGMA21(2025) journal PDF and TeX; acquired2026-09-30; exact bytes pinned by SHA256\par
\noindent\textbf{Rights:} CC BY 4.0. Authors retain copyright. \url{https://creativecommons.org/licenses/by/4.0/}. Reuse and typesetting adaptation; original language and mathematical content preserved.\par
\noindent\textbf{Difficulty (prior selection preserved):} {\CJKfont H2: 须证明任意多重切点约束的横截性、相撞标记点的交重数，并构造实现极限切触的实际曲线族。不能预设边界循环或重数公式；6.2与6.5合并为一个核心材料单元，不把各节点与尖点直接推广另计。}\par
\medskip\noindent\textbf{Editorial boundary note (not author text):} Complete original Theorem 6.2 plus its directly paired coalescence Theorem 6.5 as one existing problem unit. Retains the collision lemma and proof, intersection and closure conventions, incidence/tangency/diagonal spaces, transversality and cycle-existence argument, all multiplicity calculations, and the full local analytic coalescence construction. References to the later nodal comparison remain source references; no further singular-curve enumeration is added. The two paired identities jointly deliver the paper’s stated smooth multi-tangency enumeration method and are not counted as two separate adjacent statements.\par
\noindent\textbf{External original reference m\_rslt1:} Original Main Result 2.1, source lines 166–178, enumerative goal referenced in the section lead-in; the selected smooth-cycle and coalescence formulas are fully retained.\par
\noindent\textbf{External original reference A1FT0T0\_coll\_cycle\_ver:} Original nodal comparison equation (7.6), source lines 1727–1734, cited only in the comparison immediately before the selected smooth coalescence proof. The comparison formula itself is included verbatim; no nodal theorem is added to the selected task.\par
\medskip\hrule\medskip\noindent\textbf{Author text begins below.}\par
\setcounter{section}{3}
% Original source lines 476--1342
\section{A few remarks on intersection theory}

In this section, we summarize a few facts about intersection theory.
The ambient space $\mathsf{M}$
inside which intersection theory is done will be a product
of projective spaces. Whenever we talk about an open set,
we mean open set with the analytic topology of the projective space;
we are not talking about the Zariski topology.

Let $\alpha$ and $\beta$ be two homology classes
in $\mathsf{M}$, i.e., $\alpha, \beta \in H_*(\mathsf{M}, \mathbb{R})$.
We define $\alpha\cdot \beta$ the \emph{topological intersection} of
$\alpha$ and $\beta$ to be the unique homology class whose Poincar\'{e} dual is the
cup product of the Poincar\'{e} duals of $\alpha$ and $\beta$.

We now follow a standard abuse of notation.
Given a homology class $\alpha$, we denote its Poincar\'{e} dual (a cohomology class)
by the same letter $\alpha$.
Hence, from the point of view of cohomology,
topological intersection is simply the cup product.

The homology classes that we encounter are going to be obtained
by taking the closure of certain algebraic varieties. More precisely,
our situation is as follows: we have a set $\mathsf{S}$, which is
a smooth complex submanifold of the ambient space $\mathsf{M}$, but it is
not closed. The closure of $\mathsf{S}$ (inside $\mathsf{M}$)
need not be smooth. Nevertheless, $\overline{\mathsf{S}}$ does define a homology class.
The reason for this is as follows. The singularities of the closure are of
complex codimension one and hence, of real codimension two. Hence, by Stokes theorem
for analytic varieties
\cite[p.~33]{GH3}
integration over $\overline{\mathsf{S}}$ makes sense.
Hence, any analytic subvariety of a compact complex manifold $\mathsf{M}$, always
defines a homology class in $H_*(\mathsf{M}, \mathbb{R})$ as explained in
\cite[pp.~33 and 61]{GH3}. We often refer to~$\mathsf{S}$ as the
\emph{open part} of the cycle $\overline{\mathsf{S}}$.

We often be dealing with zero sets of sections of certain complex vector bundles.
It is a~standard fact that when the section is transverse to zero, the cycle represented
by the zero set is Poincar\'{e} dual to the top \emph{Chern class} of the vector bundle.
The top Chern class is also referred to as the \emph{Euler class} of the bundle; we
usually use the latter terminology.

\section{The collision lemma}
\label{coll_lemm}

This section contains the central lemma that we will be using throughout the paper, which is
be referred to as the collision lemma. Before stating the lemma, we need to introduce a few notations.

Let $\mathcal{D}_1$ be the space of lines in $\mathbb{P}^2$; it is the dual projective space $\check{\mathbb{P}}^2$. Define
$
M := \mathcal{D}_1 \times X_{1} \times X_{2}$,
where $X_i$ is a copy of $\mathbb{P}^2$ for $i = 1, 2$.
The pullback to $M$ of the hyperplane classes in
$X_{1}$, $X_{2}$ and $\mathcal{D}_1$ are denoted by $a_1$, $a_2$ and $y_1$, respectively. Define
\begin{align*}
X := \{ (H, q_1, q_2) \in M  \mid  q_1, q_2 \in H\}
\qquad \text{and}
\qquad
Y := \{ (H, q_1, q_2) \in X  \mid  q_1= q_2\}.
\end{align*}
\begin{lmm}[{collision lemma}]\label{CL_ver2}
The cohomology class $(a_1+a_2-y_1)$ restricted to $X$ is equal to
the Poincar\'{e} dual of $Y$ in $X$.
\end{lmm}

\begin{rem}
 Let $\Delta_{1, 2}^{\mathsf{L}}$
be the following divisor class on $M$
\begin{align}
\Delta_{1, 2}^{\mathsf{L}}& := a_1 + a_2-y_1. \label{Line_Diag_Defn}
\end{align}
The geometric content of the collision lemma
is that the intersection of $X$ with the class $\Delta_{1, 2}^{\mathsf{L}}$ is equivalent to the class
obtained by making the two points come together in $X$. \end{rem}

\begin{proof} First of all, we note that the cohomology of $X$
is generated by $a_1$, $a_2$ and $y_1$; this follows from the Leray--Hirsch theorem.
Since $Y$ is a codimension one submanifold of $X$, we conclude that the
Poincar\'{e} dual of $Y$ in $X$ is given by
\begin{align}
\textnormal{PD}_{X}[Y] = (A a_1 + B a_2 + C y_1)|_{X}, \label{k1}
\end{align}
for some numbers $A$, $B$ and $C$ that are to be determined.
To prove the lemma, it suffices to show that $A = 1$, $ B = 1$ and $C = -1$.
This, we show by computing three different intersection numbers.

First of all, we note that the Poincar\'{e} dual of $X$ in $M$ is given by
\begin{align}
\textnormal{PD}_{M}[X]& = (y_1+a_1)\cdot(y_1+a_2). \label{k2}
\end{align}
To see why this is so, first define $Z_1$ as
$
Z_1 := \{ (H, q_1) \in \mathcal{D}_1 \times X_1  \mid  q_1 \in H\}$.
To prove \eqref{k2}, it suffices to show that $[Z_1] = (y_1 + a_1)$.
To justify this, note that
$[Z_1] = n y_1 + m a_1$ for some numbers $n$ and $m$. Next, we note that
$n = [Z_1]\cdot y_1 a_1^2$ and $m = [Z_1]\cdot y_1^2 a_1$.
Geometrically, $[Z_1]\cdot y_1^2 a_1$ is the number of lines passing through
two points \big(which corresponds to the factor~$y_1^2$\big) and a marked point on the line
that intersects another generic line (which corresponds to the factor $a_1$). This is
clearly $1$. Similarly, $[Z_1]\cdot y_1 a_1^2$ is geometrically the number of lines
through two points (the first point corresponds to the factor $a_1^2$ and the second point
corresponds to the factor $y_1$). This number is also equal to $1$. Hence,
$[Z_1]= y_1 + a_1$, thereby proving \eqref{k2}.

Next, we note that $Y$ can also be described as a codimension three submanifold of $M$
in the following way
\begin{align*}
Y:= \{ (H, q_1, q_2) \in M\mid q_1 \in H,\,  q_1 = q_2\}.
\end{align*}
Hence, the Poincar\'{e} dual of $Y$ in $M$ is given by
$
\textnormal{PD}_{M}[Y] = (y_1+a_1)\cdot [\Delta_{12}]$,
where $\Delta_{12}$ is the subspace of points $(H_1, q_1, q_2)$ in $M$, such that
$q_1 = q_2$.
Using the standard fact that the class of the diagonal
$[\Delta_{12}]$ is equal to $\big(a_1^2 + a_1 a_2 + a_2^2\big)$,
we conclude that
\begin{align}
\textnormal{PD}_{M}[Y]& = (y_1+a_1)\cdot \big(a_1^2 + a_1 a_2 + a_2^2\big). \label{k3}
\end{align}
Let $\mu$ be a class in $M$ of degree three. By equations \eqref{k1}, \eqref{k2} and \eqref{k3},
we conclude that
\begin{align*}
(y_1+a_1)\cdot (a_1^2 + a_1 a_2 + a_2^2) \cdot \mu & = (A a_1 + B a_2 + C y_1)\cdot (y_1+a_1)\cdot(y_1+a_2) \cdot \mu.
\end{align*}
By suitably choosing $\mu$, we can determine $A$, $B$ and $C$.
Choosing $\mu := a_1 a_2^2$, $a_2a_1^2$ and $a_2 y_1^2$, we have
$A+C = 0$, $B+C = 0$ and $A = 1$, respectively. This precisely implies that
$A = B = 1$ and $C = -1$.
\end{proof}

\section{Counting smooth curves with multiple tangencies}\label{smooth_tangencies}
\label{smooth_curves_enum}
In this section, we use the collision lemma to derive our Main Result~\ref{m_rslt1}.
Before we get into the details, a brief outline of our method is described.

Let $\mathcal{D}_d$
denote the space of degree $d$ curves in $\mathbb{P}^2$. This is a complex projective space of dimension
\begin{equation}\label{eb}
\delta_d :=  \frac{d(d+3)}{2}.
\end{equation}
We now try to answer the following question:
{\it How many degree $d$ curves are there in $\mathbb{P}^2$ passing through
$\delta_d-1$ generic points and that are tangent to a given line?}

We approach this problem in the following way: First, consider the space of curves with two marked points $x_1$ and $x_2$, such that
both the points lie on the curve and the line. Now
impose the condition that $x_1$ becomes equal to $x_2$; this is precisely where
we use the collision lemma.
Once we impose the condition that the points $x_1$ and $x_2$
have become equal,
the curve becomes tangent to this line.
We now implement this idea precisely.
Let $X_1$ denote a copy of the projective plane. Define
the incidence variety
\begin{equation}\label{f1}
\mathrm{I}_d := \{ (H_d, x_1) \in  \mathcal{D}_d \times X_1 \mid  x_1 \in H_d\}.
\end{equation}
Elements of the incidence variety consists of
degree $d$-curve $H_d$ and a marked point $x_1$ that lies on this curve.
Let $a_1$ and $y_d$ denote the divisor classes on $\mathcal{D}_d \times X_1$ obtained by pulling back the
hyperplane classes on $X_1$ and $\mathcal{D}_d$, respectively. Then we have
\begin{equation}\label{incidence}
[\mathrm{I}_d]  =  y_d + d a_1.
\end{equation}
Indeed, this follows immediately from the fact that the restrictions of \eqref{incidence} to both
$\{{\rm point}\}\times \mathbb{CP}^2$ and $\mathcal{D}_d \times \{{\rm point}\}$ are valid.

Next, we study subspaces of $\mathcal{D}_1 \times \mathcal{D}_d \times X_1$.
First, define
$\mathrm{I}_{\mathsf{L}}$ (respectively, $\mathrm{I}_{\mathcal{C}}$)
to be the subspace of $\mathcal{D}_1 \times \mathcal{D}_d \times X_1$ consisting of all
$(H_1, H_d, x_1)  \in \mathcal{D}_1 \times \mathcal{D}_d \times X_1$
such that $x_1 \in H_1$ (respectively, $x_1 \in H_d$).
Next, define
\begin{equation}
\label{et}
\mathsf{T}_0
\end{equation}
to be the subspace
of $\mathcal{D}_1 \times \mathcal{D}_d \times X_1$ consisting of all
$(H_1, H_d, x_1)  \in \mathcal{D}_1 \times \mathcal{D}_d \times X_1$ such that~$H_1$ and~$H_d$ intersect transversally at $x_1$.
Note that the closure $\overline{\mathsf{T}}_0$ consists of
$(H_1, H_d, x_1)  \in \mathcal{D}_1 \times \mathcal{D}_d \times X_1$
such that $x_1 \in H_1\bigcap H_d$.

For making the notation easier to read, we denote the homology class
represented by the closure by the notation $[\mathsf{T}_0]$
\big(as opposed to more cumbersome $\big[\overline{\mathsf{T}}_0\big]$\big).

\begin{lmm}\label{Theo_tzr}
The class $[\mathsf{T}_0] \in H_{2(\delta_d +2)}(\mathcal{D}_1 \times \mathcal{D}_d \times X_1, {\mathbb R})$
represented by the cycle $\overline{\mathsf{T}}_0 \subset \mathcal{D}_1 \times \mathcal{D}_d \times X_1$
$($see \eqref{eb}, \eqref{et} and \eqref{eb}$)$ is the following:
\begin{equation}\label{T0_cycle_expression_base}
[\mathsf{T}_0]  =  [\mathrm{I}_{\mathsf{L}}] \cdot [\mathrm{I}_{\mathcal{C}}],
\end{equation}
where $[\mathrm{I}_{\mathsf{L}}]$ and $[\mathrm{I}_{\mathcal{C}}]$ are the homology classes of
$\mathrm{I}_{\mathsf{L}}$ and $\mathrm{I}_{\mathcal{C}}$, respectively.
\end{lmm}

\begin{proof}
Consider the divisor $\mathrm{I}_1  \subset \mathcal{D}_1\times X_1$ in \eqref{f1}.
Let
$
f \colon  \mathrm{I}_1\times \mathcal{D}_d  \hookrightarrow  \mathcal{D}_1 \times \mathcal{D}_d \times X_1$
be the natural map defined by $((H_1, x_1), H_d) \longmapsto (H_1, H_d, x_1)$, where
$(H_1, x_1) \in \mathrm{I}_1$ and $H_d \in \mathcal{D}_d$. Let
\begin{equation}\label{f3}
p \colon\ \mathcal{D}_1 \times \mathcal{D}_d \times X_1 \longrightarrow
\mathcal{D}_d \times X_1
\end{equation}
be the natural projection. Note that
$
\overline{\mathsf{T}}_0 \subset \mathrm{I}_1\times \mathcal{D}_d
$
(see \eqref{et}) and $\overline{\mathsf{T}}_0$ is a divisor.

To prove the lemma, it suffices to show the following:
\begin{equation}\label{f4}
f^*p^* {\mathcal O}_{\mathcal{D}_d \times X_1}(\mathrm{I}_d) =
{\mathcal O}_{\mathrm{I}_1\times \mathcal{D}_d}\bigl(\overline{\mathsf{T}}_0\bigr).
\end{equation}

We now describe three subvarieties $S_1$, $S_2$ and $S_3$ of $\mathrm{I}_1\times \mathcal{D}_d$. Fix
a point $z_0 \in \mathrm{I}_1$ and define%
\begin{equation}\label{s1}
S_1 := \{z_0\}\times\mathcal{D}_d \subset \mathrm{I}_1\times \mathcal{D}_d.
\end{equation}
Fix curves $(H_1, H_d) \in \mathcal{D}_1 \times \mathcal{D}_d$ in $X_1$, and we have
\begin{equation}\label{s2}
S_2 := \{(H_1, H_d, x_1) \in \mathcal{D}_1 \times \mathcal{D}_d \times X_1
  \mid  x_1 \in H_1\} \subset \mathrm{I}_1\times \mathcal{D}_d.
\end{equation}
So $S_2$ is identified with the line $H_1$. Fix a point $x_0 \in X_1$ and also an
element $H_d \in \mathcal{D}_d$. Now define
\begin{equation}\label{s3}
S_3 := \{(H, H_d, x_0) \in \mathcal{D}_1 \times \mathcal{D}_d \times X_1  \mid
x_0 \in H\} \subset \mathrm{I}_1\times \mathcal{D}_d.
\end{equation}
Consequently, $S_3$ is identified with the pencil of lines in $X_1$ containing $x_0$.

Given two cohomology classes ${\mathcal L}_1, {\mathcal L}_2  \in H^2(\mathrm{I}_1\times \mathcal{D}_d,
{\mathbb R})$, to show that
${\mathcal L}_1 = {\mathcal L}_2$, it is enough~to prove that
\smash{$
{\mathcal L}_1\big\vert_{S_j} = {\mathcal L}_2\big\vert_{S_j}
$}
for $j = 1, 2, 3$ (see \eqref{s1}, \eqref{s2}, \eqref{s3}).

In view of the above criterion, it is now straightforward to prove \eqref{f4} using \eqref{incidence}.
\end{proof}

Note that using \eqref{incidence}, we can rewrite equation \eqref{T0_cycle_expression_base} as
\begin{equation}\label{T0_cycle_expression_base_ag}
[\mathsf{T}_0]  =  (y_1 + a_1)\cdot(y_d + d a_1),
\end{equation}
where $\cdot$ denotes topological intersection. Generalizing the notion of
$\mathsf{T}_0$, given any nonnegative integer $k$, we define $\mathsf{T}_k$
to be the subspace of $\mathcal{D}_1 \times \mathcal{D}_d\times X_1$
consisting of all points $(H_1, H_d, x_1)$ such that
\begin{itemize}\itemsep=0pt
\item The points $x_1$ is a smooth point of the curve $H_d$.
\item The line $H_1$
intersects the curve $H_d$ at the points $x_1$
with the order of tangency precisely equal to $k$.
\end{itemize}
Notice that as per its definition, $\mathsf{T}_1$ is not a subset of $\mathsf{T}_0$, but a subset of the
closure $\overline{\mathsf{T}}_0$.

Next, generalizing $\mathcal{D}_1 \times \mathcal{D}_d \times X_1$, define
$
\mathsf{M}_n :=  \mathcal{D}_1 \times \mathcal{D}_d \times (X_1\times \cdots\times X_n)$,
where $X_j$, $1 \leq j \leq n$, is a copy of $\mathbb{P}^2$.
Let $\mathsf{S}_1, \mathsf{S}_2, \dots, \mathsf{S}_n$ be subvarieties of
$\mathcal{D}_1 \times \mathcal{D}_d \times \mathbb{P}^2$.
Denote $\mathsf{S}_1 \mathsf{S}_2 \dots \mathsf{S}_n$ by $\mathsf{S}$. Then,
$\mathsf{S} \subset \mathsf{M}_n$
consists of
$(H_1, H_d, x_1,\dots, x_n)$ such that
\begin{itemize}\itemsep=0pt
\item $(H_1, H_d, x_i) \in \mathsf{S}_i$ for all $i=1$ to $n$.
\item The points $x_1,\dots, x_n$ are all distinct.
\end{itemize}
As an example, consider the set $\mathsf{T}_1 \mathsf{T}_2$. This comprises of the set of
curves along with two \textit{distinct} marked points on a line, where the curve is tangent to first order to
the line at the first marked point and is tangent to second order to the line at the
second marked point.
Similarly, $\overline{\mathsf{T}}_1 \overline{\mathsf{T}}_2$~denotes a slightly bigger space, where the curve is
at least as degenerate as being tangent to the line to first order at the first marked
point and is at least
as degenerate as being tangent to the line to second order at the second marked point.
The curve could be tangent to second order at the first marked point. The curve could even have a nodal
point at the first marked point, since these both lie in the closure $\mathsf{T}_1$.
However, the two marked points have to be distinct.
In particular, $\overline{\mathsf{T}}_1 \overline{\mathsf{T}}_2$
is not the set-theoretic intersection of
$\overline{\mathsf{T}}_1$ and $\overline{\mathsf{T}}_2$, since the latter
includes the locus where the two marked points are equal.
Finally, in the space
$\overline{\mathsf{T}_1 \mathsf{T}}_2$, the two marked points need not be distinct; this denotes the
closure of the space $\mathsf{T}_1 \mathsf{T}_2$ and it includes the locus where the two marked points coincide.

We denote the homology class defined by the closure of $\mathsf{S}$
by the notation $\big[\mathsf{S}\big]$ as opposed to the more cumbersome $[\overline{\mathsf{S}}]$; this makes
some of the computations and statements easier to read.

Finally, let
\begin{equation}\label{ep}
\pi \colon\ \mathsf{M}_{n+1} \longrightarrow \mathsf{M}_n
\end{equation}
be the projection that sends
any $(H_1, H_d, x_1, \dots, x_{n}, x_{n+1})$ to $(H_1, H_d, x_1, \dots, x_{n})$. Let
\begin{equation}\label{ep1}
\pi_{n+1} \colon\  \mathsf{M}_{n+1} \longrightarrow \mathsf{M}_1
\end{equation}
be the projection that sends any $(H_1, H_d, x_1, \dots, x_{n}, x_{n+1})$ to
$(H_1, H_d, x_{n+1})$. For any $1 \leq i \leq n$, let
$
\Delta_{i, {n+1}}  \subset  \mathsf{M}_{n+1}
$
be the locus of all $(H_1, H_d, x_1, \dots, x_{n}, x_{n+1})$ such that $x_i = x_{n+1}$.
We are now ready to state the main results to enumerate smooth curves with tangencies.

\begin{thm}\label{theorem_for_many_Tks}
Let $n$ be a positive integer and $k_1, k_2,\dots, k_{n}$ nonnegative integers
with
$
k :=  k_1+k_2+\dots +k_n$.
Then the following equality of elements of $H_{2(\delta_d+2- k)}(\mathsf{M}_{n+1}, {\mathbb R})$ holds:
\begin{equation}\label{many_Tks}
\pi^{\ast}[\mathsf{T}_{k_1}\dots \mathsf{T}_{k_n}]\cdot
\pi_{n+1}^{\ast}[\mathsf{T}_0] =
[\mathsf{T}_{k_1}\dots \mathsf{T}_{k_n} \mathsf{T}_{0}]
 + \sum_{i=1}^{n}(k_i+1) \pi^{\ast}[\mathsf{T}_{k_1}\dots \mathsf{T}_{k_n}] \cdot [\Delta_{i, {n+1}}]
\end{equation}
$($see \eqref{ep}, \eqref{ep1}$)$ provided $d > k+n$.
\end{thm}

Note that implicit in Theorem \ref{theorem_for_many_Tks} is the assertion that
the closure of $\mathsf{T}_{k_1}\dots \mathsf{T}_{k_n}$ in
$\mathsf{M}_n$ and the closure of
$\mathsf{T}_{k_1}\dots \mathsf{T}_{k_n} \mathsf{T}_{0}$ in $\mathsf{M}_{n+1}$
actually define homology classes. We will justify that assertion as well.
In order to prove Theorem \ref{theorem_for_many_Tks}, we first prove an intermediate statement
which is interesting in its own right.
\begin{prp}
\label{prp_smth_mfld}
Let $n$ be a positive integer and $k_1, k_2,\dots, k_{n}$ nonnegative integers
with
$ k :=  k_1+k_2+\dots +k_n$.
Then $\overline{\mathsf{T}}_{k_1} \overline{\mathsf{T}}_{k_2}\dots \overline{\mathsf{T}}_{k_n}$
is a smooth submanifold of $\mathsf{M}_{n}$, provided
$d \geq k+n$.
\end{prp}

\begin{rem}
 We are \emph{not} claiming that
\smash{$\overline{\mathsf{T}_{k_1}\mathsf{T}_{k_2}\dots \mathsf{T}}_{k_n}$}
is a smooth manifold. Notice the difference between
\smash{$\overline{\mathsf{T}}_{k_1} \overline{\mathsf{T}}_{k_2}\dots \overline{\mathsf{T}}_{k_n}$}
and \smash{$\overline{\mathsf{T}_{k_1}\mathsf{T}_{k_2}\dots \mathsf{T}}_{k_n}$}; in the former space
all the marked points are distinct, while in the latter space, that is not necessarily true.
\end{rem}

\begin{proof}
We start by proving the proposition for $n = 1$
and $k_1 = 0$, i.e., we show that $\overline{\mathsf{T}}_0$ is a~smooth
submanifold of $\mathsf{M}_1$.

Let $\mathcal{F}_d$ be the space of
polynomials in two variables of degree at most $d$.
This is a vector space of dimension $\frac{d(d+3)}{2} + 1$, because
an element of $\mathcal{F}_d$ can be viewed as
\[
 f(x, y) := f_{00} + f_{10}x + f_{01} y + \frac{f_{20}}{2} x^2 + f_{11} xy + \frac{f_{02}}{2} y^2 +
\dots + \frac{f_{0d}}{d!} y^d,
\]
and hence $f$ can be identified with the vector
$(f_{00}, f_{10}, f_{01},\dots, f_{0d})$.
Let $\mathcal{F}_d^{+}$ denote the space of nonzero polynomials of degree $d$.
We note that the projectivization of $\mathcal{F}_d^{+}$ is $\mathcal{D}_d$.

Consider the projection map
\[ \pi_{+}\colon\ \mathcal{F}_1^{+} \times \mathcal{F}_d^{+} \times X_1 \longrightarrow
\mathcal{D}_1 \times \mathcal{D}_d \times X_1.  \]
We note that the projection map is a submersion. Hence, it suffices to show that
\smash{$\widetilde{\overline{\mathsf{T}}}_0:=
\pi_{+}^{-1}\bigl(\overline{\mathsf{T}}_0\bigr)$} is a smooth submanifold of
$\mathcal{F}_1^{+} \times \mathcal{F}_d^{+} \times X_1$. We prove that now.

Let \smash{$(g, f, p_1) \in \widetilde{\overline{\mathsf{T}}}_0$}.
By choosing a chart around the point $p_1$, we can identify an open neighbourhood of
$p_1$ (in $X_1$) by $\mathbb{C}^2$. Hence, in order to show that
\smash{$\widetilde{\overline{\mathsf{T}}}_0$} is a smooth
submanifold of~${\mathcal{F}_1^{+} \times \mathcal{F}_d^{+} \times X_1}$, it suffices to show that
$(0,0)$ is a regular value of the map
\[\varphi \colon\ \mathcal{F}_1^{+} \times \mathcal{F}_d^{+} \times \mathbb{C}^2 \longrightarrow
\mathbb{C}^2,
\qquad \text{given by}
\quad \varphi(g, f, (x, y)) := (g(x, y), f(x, y)).\]
Let us prove that now. Assume that $\varphi(g,f,(a,b))=0$.
We need to show that the differential of~$\varphi$,
evaluated at $(g, f, (a, b))$ is surjective.
In order to prove that, consider the two curves
\begin{align*}
\gamma_1, \gamma_2&\colon\ (-\varepsilon, \varepsilon)\longrightarrow
\mathcal{F}_1^{+} \times \mathcal{F}_d^{+} \times \mathbb{C}^2,
\qquad \textnormal{given by} \\
\gamma_1(t)&:= (g + t\eta_1, f, (a, b)) \qquad \textnormal{and}
\qquad \gamma_2(t):= (g, f+t\eta_2, (a, b)),
\end{align*}
where
$\eta_1$ and $\eta_2$ are as yet, unspecified elements of
$\mathcal{F}_1^{+}$ and $\mathcal{F}_d^{+}$.
We now note that
\[\{{\rm d}\varphi|_{(g,f(a,b))}\}(\gamma_1^{\prime}(0)) = (\eta_1(a,b), 0)
\qquad \textnormal{and} \qquad \{{\rm d}\varphi|_{(g,f,(a,b))}\}(\gamma_2^{\prime}(0)) = (0, \eta_2(a,b)).\]
Hence, to prove that the differential is surjective, we simply need to produce
$\eta_1 \in \mathcal{F}_1^{+}$ and $\eta_2 \in \mathcal{F}_d^{+}$
such that $\eta_1(a,b) \neq 0$ and $\eta_2(a,b) \neq 0$.
That is easily achieved: we simply define both of them to be the constant functions
taking the value $1$.

Before proceeding further, we make a couple of simplifications that make the
subsequent proofs easier.
We showed that if $\varphi(g, f, (a, b)) = 0$, then the
differential of $\varphi$ is surjective. We claim that by making a suitable change of
coordinates, we can always assume that~$(a,b)$ is the origin and the line is the $x$-axis.
To see why this is so,
assume that the line is given by~${g_{00} + g_{10}x + g_{01}y = 0}$.
Assuming that $g_{01}\neq 0$, define the new coordinates $X$
and $Y$ by $X := x-a$ and $Y := g_{00} + g_{10}x + g_{01}y$.
If $g_{01} =0$, then define $X:= y-b$ and $Y:= g_{00} + g_{10}x + g_{01}y$.
Define
\[F(X,Y):= f(x(X,Y), y(X,Y)).\]
This is the polynomial $f$
written in the new coordinates $X$ and $Y$. Define
\smash{$F_{ij} := \frac{\partial^{i+j}F(X,Y)}{\partial X^i \partial Y^j}\big|_{(0,0)}$}.
In these new coordinates, the point under consideration is the origin
and the line is the $X$-axis. Furthermore, the coefficients of the polynomial
are given by $\frac{F_{ij}}{i!j!}$.
Hence, what we have shown is that
\smash{$\widetilde{\overline{\mathsf{T}}}_0$} is a fibre bundle over the incidence variety
$J$ (where $J$ is defined to be the subset of~${\mathcal{F}_1^{+}\times X_1}$ where the point lies on the line)
and the fibre over $(g, p) \in J$ can be identified with~\smash{$\big(\overline{\mathsf{T}}_0\big)_{\mathsf{Aff}}$}, where
$\big(\overline{\mathsf{T}}_0\big)_{\mathsf{Aff}} := \{ f \in \mathcal{F}_d^+\colon f_{00}=0 \}$.
The map $f$ going to $F$ is a trivialization of this fibre bundle.
Henceforth, we set the line to be the $x$-axis and the point to be the origin; this makes the
calculations simpler.

We now show that $\overline{\mathsf{T}}_{k_1}$ is a smooth
submanifold of $\mathsf{M}_1$ for all $k_1$. We use induction on $k_1$.
Assume that we have proved the assertion till $k_1-1$. Hence
$\overline{\mathsf{T}}_{k_1-1}$ is a smooth
submanifold of~$\mathsf{M}_1$. Hence,
\smash{$\widetilde{\overline{\mathsf{T}}}_{k_1-1}:=
\pi_{+}^{-1}\big(\overline{\mathsf{T}}_{k_1-1}\big)$} is a smooth submanifold of
$\mathcal{F}_1^{+} \times \mathcal{F}_d^{+} \times X_1$.

We now show that \smash{$\widetilde{\overline{\mathsf{T}}}_{k_1}$} is a smooth submanifold
of \smash{$\widetilde{\overline{\mathsf{T}}}_{k_1-1}$}.
Define \smash{$\big(\overline{\mathsf{T}}_{k_1-1}\big)_{\mathsf{Aff}}$} as follows:
\[ \bigl(\overline{\mathsf{T}}_{k_1-1}\bigr)_{\mathsf{Aff}} :=
\bigl\{ f \in \mathcal{F}_d^{+}\mid f_{00}, f_{10} \dots, f_{k_1-1, 0} =0\bigr\}. \]
We note that
$\widetilde{\overline{\mathsf{T}}}_{k_1-1}$
is a fibre bundle over $\mathcal{F}_1^+\times X_1$
whose fibres can be identified with
\smash{$\bigl(\overline{\mathsf{T}}_{k_1-1}\bigr)_{\mathsf{Aff}}$}.
In order to show that
\smash{$\widetilde{\overline{\mathsf{T}}}_{k_1}$} is a smooth submanifold of
$\mathcal{F}_1^+\times \mathcal{F}_d^+\times X_1$, it suffices to show that
zero is a regular value of the map
\smash{$\varphi\colon \bigl(\overline{\mathsf{T}}_{k_1-1}\bigr)_{\mathsf{Aff}} \longrightarrow
\mathbb{C}$}, given by
$\varphi(f):= f_{k_1, 0}$.
Suppose
$f \in \bigl(\overline{\mathsf{T}}_{k_1-1}\bigr)_{\mathsf{Aff}}$.
Let \smash{$\gamma\colon(-\varepsilon, \varepsilon)\longrightarrow
\bigl(\overline{\mathsf{T}}_{k_1-1}\bigr)_{\mathsf{Aff}}$}
be a curve, given by
$\gamma(t):= f+t\eta$,
where $\eta$ is as yet an unspecified polynomial. We now note that
$\{d \varphi\}(\gamma^{\prime}(0)) = \eta_{k_1, 0}$. Now choose $\eta$ as follows
$\eta(x, y) := x^{k_1}$.
Since
$f \in \bigl(\overline{\mathsf{T}}_{k_1-1}\bigr)_{\mathsf{Aff}}$,
$f + t \eta$ also belongs to $\bigl(\overline{\mathsf{T}}_{k_1-1}\bigr)_{\mathsf{Aff}}$
for all $t$ nonzero but small.
Furthermore, $\eta_{k_1, 0} \neq 0$. This proves the claim.

Next, for multiple points, we now use induction on
$n$. If $k_n \geq 1$, define
\[\bigl(\overline{\mathsf{T}}_{k_1} \overline{\mathsf{T}}_{k_2}
\dots \overline{\mathsf{T}}_{k_{n-1}} \overline{\mathsf{T}}_{k_{n}-1}\bigr)_{\mathsf{Aff}}
\]
to be the following subset of $\mathcal{F}_d^{+} \times \mathbb{C}^{n-1}$: it is
the collection of all $(f, \mathbf{a}_1, \dots, \mathbf{a}_{n-1})$, such that
\begin{itemize}\itemsep=0pt
\item The numbers
$\mathbf{a}_1, \mathbf{a}_2, \dots \mathbf{a}_{n-1}$ are all distinct
from each other and different from zero.
\item All the derivatives of $f$ with respect to $x$
at $\mathbf{a}_i$ up to order $k_i$ are zero, for $i=1$ to $n-1$.
\item All the derivatives of $f$ with respect to $x$
at $(0, 0)$ up to order $k_n-1$ are zero.
\end{itemize}
It is also convenient to define
$\bigl(\overline{\mathsf{T}}_{k_1} \overline{\mathsf{T}}_{k_2}
\dots \overline{\mathsf{T}}_{k_{n-1}} \overline{\mathsf{T}}_{-1}\bigr)_{\mathsf{Aff}}$
as the following subset of $\mathcal{F}_d^{+} \times \mathbb{C}^{n-1}$: it is
the collection of all $(f, \mathbf{a}_1, \dots, \mathbf{a}_{n-1})$, such that
\begin{itemize}\itemsep=0pt
\item The numbers
$\mathbf{a}_1, \mathbf{a}_2, \dots , \mathbf{a}_{n-1}$ are all distinct
from each other and different from zero.
\item All the derivatives of $f$ with respect to $x$
at $\mathbf{a}_i$ up to order $k_i$ are zero, for $i=1$ to $n-1$.
\end{itemize}
Arguing as before,
it suffices to show that for all $k_n \geq 0$, zero is a regular value of the map
\[\varphi \colon\ \bigl(\overline{\mathsf{T}}_{k_1} \overline{\mathsf{T}}_{k_2}
\dots \overline{\mathsf{T}}_{k_{n-1}} \overline{\mathsf{T}}_{k_{n}-1}\bigr)_{\mathsf{Aff}}
 \longrightarrow \mathbb{C} \qquad \text{given by}
\quad \varphi(f, x_1, \dots, x_{n-1}) := f_{k_n, 0}. \]
Let us prove that now. Suppose
\[(f,\mathbf{a}_1, \dots, \mathbf{a}_{n-1})
\in \bigl(\overline{\mathsf{T}}_{k_1} \overline{\mathsf{T}}_{k_2}
\dots \overline{\mathsf{T}}_{k_{n-1}} \overline{\mathsf{T}}_{k_{n}-1}\bigr)_{\mathsf{Aff}}.\]
Let $\gamma\colon(-\varepsilon, \varepsilon)\longrightarrow
\bigl(\overline{\mathsf{T}}_{k_1} \overline{\mathsf{T}}_{k_2}
\dots \overline{\mathsf{T}}_{k_{n-1}} \overline{\mathsf{T}}_{k_{n}-1}\bigr)_{\mathsf{Aff}}$
be a curve, given by
$\gamma(t) := (f+t\eta, \mathbf{a}_1, \dots, \mathbf{a}_{n-1})$
where $\eta$ is as yet an unspecified polynomial. We now note that
$\{{\rm d} \varphi\}(\gamma^{\prime}(0)) = \eta_{k_{n}, 0}$. Now choose~$\eta$ as follows
\[ \eta(x, y) := (x-\mathbf{a}_1)^{k_1+1} \cdots (x-\mathbf{a}_{n-1})^{k_{n-1}+1} (x-0)^{k_n}.\]
Since $(f, \mathbf{a}_1, \dots, \mathbf{a}_{n-1})$
belongs to $\bigl(\overline{\mathsf{T}}_{k_1} \overline{\mathsf{T}}_{k_2}
\dots \overline{\mathsf{T}}_{k_{n-1}} \overline{\mathsf{T}}_{k_{n}-1}\bigr)_{\mathsf{Aff}}$,
$\gamma(t)$
also lies there for all $t$.
Furthermore, $\eta_{k_n, 0} \neq 0$, because the $\mathbf{a}_i$ are all different from zero.
This proves the proposition.
\end{proof}

We are now ready to prove Theorem \ref{theorem_for_many_Tks}.

\begin{proof}[Proof of Theorem \ref{theorem_for_many_Tks}]
We first show that \eqref{many_Tks} is valid on the set-theoretic level. Consider the first term
on the left-hand side, namely
$\pi^{\ast}[\mathsf{T}_{k_1}\dots \mathsf{T}_{k_n}]$.
It is represented by the closure of the following space: a line, a curve and
$n+1$ distinct points $(x_1,\dots, x_{n+1})$, such that the
curve is tangent to the line at the points
$(x_1,\dots, x_n)$ to orders $k_1,\dots, k_n$
respectively; the last point~$x_{n+1}$ is free (it does not have to lie on either the line or the curve).

Now consider second factor on the left-hand side of \eqref{many_Tks},
namely $\pi_{n+1}^{\ast}[\mathsf{T}_0]$.
This is simply represented by the following space: a
line, a curve and $n+1$ points $(x_1,\dots, x_{n+1})$,
such that the points $(x_1,\dots, x_n)$ are free, while the last point $x_{n+1}$ lies on the line
and the curve.

Consider the set-theoretic intersection of the above two spaces. There are two possibilities. The first
possibility is that the point $x_{n+1}$ is distinct from all the other points
$(x_1,\dots, x_n)$.
The closure of that space represents the first term on the right-hand side of
\eqref{many_Tks}. But there is another possibility. The point $x_{n+1}$ could be equal to one of the
$x_i$
(for $i \in \{1,\dots, n\}$).
That precisely gives us the second term on the right-hand side of \eqref{many_Tks}.

To see that equation \eqref{many_Tks} is valid on the level of homology,
we need to do the following.
To justify the first term on the right-hand side of \eqref{many_Tks},
we need to
show that the intersections are transverse;
this follows from the proof of Proposition \ref{prp_smth_mfld}.
To justify the second term,
we need to justify the multiplicity of the intersection.

Consider the situation of
$x_1$ coinciding with $x_{n+1}$. We take a chart that sends the point
$x_{n+1}$ to be the origin and sets the line to be the $x$-axis.
The situation now is that we have a curve $f$ that is tangent to the $x$-axis to order $k_1$
at the origin. We are now going to
study the multiplicity with which the evaluation map vanishes at the origin.
Hence, $f$ is such that $f_{00}, f_{10}, \dots , f_{k_1 0}$ all vanish. It is given by
\[
f(x, y)  =  \frac{f_{k_1+1, 0}}{(k_1+1)!} x^{k_1+1}
+\frac{f_{k_1+2, 0}}{(k_1+2)!} x^{k_1+2} + \dots +
\frac{f_{d, 0}}{d !} x^{d}  + y \mathcal{R}(x, y).
\]
Now consider the evaluation map
\[
\varphi(f, x) :=  f(x, 0)  = \frac{f_{k_1+1, 0}}{(k_1+1)!} x^{k_1+1}
+\frac{f_{k_1+2, 0}}{(k_1+2)!} x^{k_1+2} + \dots +
\frac{f_{d, 0}}{d !} x^{d}.
\]
The order of vanishing of $\varphi$ is clearly $k_1+1$, provided
$f_{k_1+1, 0} \neq 0$
(the values of $f_{k_1+2, 0}, f_{k_1+3, 0},\allowbreak \dots, f_{d,0}$ are not relevant for the order of vanishing
in a neighbourhood of the origin if $f_{k_1+1, 0}$ is non-zero).

The assumption $f_{k_1+1, 0} \neq 0$ is valid, because
to compute the order of vanishing, we will be intersecting with
cycles that correspond to constraints being generic. Hence, the order of vanishing is $k_1+1$.
This proves \eqref{many_Tks} on the level of homology.
\end{proof}

We are now ready to prove our next result.

\begin{thm}
\label{Tk1Tk2_etc_alt_way}
Let $n$ be a positive integer, $k_1, k_2,\dots, k_{n-1}$ nonnegative integers
and $k_n$ a~positive integer.
Define
$
k :=  k_1+k_2+\dots +k_n$.
Then the following equality of elements of $H_{2(\delta_d-k)}(\mathsf{M}_{n+1}, {\mathbb R})$ holds:
\begin{align}
[\mathsf{T}_{k_1} \dots \mathsf{T}_{k_{n-1}} \mathsf{T}_{k_n-1} \mathsf{T}_0] \cdot \bigl[\Delta_{n, {n+1}}^{\mathsf{L}}\bigr] & =
\pi^*[\mathsf{T}_{k_1} \dots \mathsf{T}_{k_n}] \cdot [\Delta_{n, n+1}],
\label{Tmu_T0_nu_alt_way_ag}
\end{align}
provided $d > k+n-1$.
\end{thm}

Before we prove Theorem \ref{Tk1Tk2_etc_alt_way}, a few things are explained.
Consider the special case of this theorem, when $n = 1$ and $k_1 = 1$.
In this case, \eqref{Tmu_T0_nu_alt_way_ag} simplifies to
\begin{align}
[\mathsf{T}_0 \mathsf{T}_0] \cdot \bigl[\Delta^{\mathsf{L}}_{12}\bigr] &  =  \pi^*[\mathsf{T}_1]\cdot [\Delta_{12}].
\label{T0T0_collides_T1_aagg}
\end{align}
If we draw an analogy with equation \eqref{A1FT0T0_coll_cycle_ver},
it might seem that the right-hand side
of equation~\eqref{T0T0_collides_T1_aagg}
has a missing term, namely
a term that corresponds to nodal curves lying
on a line. However, that is not the case; there are no missing terms.
The term that seems to be missing is actually present: it is present
inside the \emph{closure} $\overline{\mathsf{T}}_1$. However, since this locus
is one codimension higher, when we intersect
equation \eqref{T0T0_collides_T1_aagg} with a class of complementary dimension, we do not get
any contribution. We explain this more precisely. Define
$\mu := y_1^{2} y_d^{\delta_d-1}$. Now intersect both sides of
equation \eqref{T0T0_collides_T1_aagg} with $\mu$. That gives us
$[\mathsf{T}_0 \mathsf{T}_0] \cdot \bigl[\Delta^{\mathsf{L}}_{12}\bigr]\cdot \mu  =
[\mathsf{T}_1]\cdot \mu$.
The term that one might be worried that one has missed out, namely nodal curves
with the node lying on the line, giving empty intersection with $\mu$; this is
because we are making the curve pass through $\delta_d-1$ generic points \big(since we
are intersecting with $y^{\delta_d-1}$\big).

In contrast, look at equation \eqref{A1FT0T0_coll_cycle_ver}, namely
\begin{align*}
\bigl[\mathsf{A}_1^{\mathsf{F}}\mathsf{T}_0 \mathsf{T}_0\bigr] \cdot \bigl[\Delta^{\mathsf{L}}_{12}\bigr] &  =  \pi^*\bigl[\mathsf{A}_1^{\mathsf{F}}\mathsf{T}_1\bigr]\cdot [\Delta_{12}] + 2\pi^*\bigl[\mathsf{A}_1^{\mathsf{L}}\bigr]\cdot [\Delta_{12}]\cdot
\bigl[\Delta_{2}^1\bigr].
\end{align*}
The first geometric fact we note that $\mathsf{A}_1^{\mathsf{L}}$ is \emph{not} a subset of the
closure \smash{$\overline{\mathsf{A}_1^{\mathsf{F}}\mathsf{T}}_1$}.
What is true is that $\mathsf{A}_1^{\mathsf{F}}\mathsf{A}_1^{\mathsf{L}}$
is a subset of the closure (which for dimensional reasons, does not contribute when
we intersect with a class of complimentary dimension). In order to extract
numbers, define~$\mu := y_1^2 y_d^{\delta_d-2}$.
Intersecting both sides of equation \eqref{A1FT0T0_coll_cycle_ver}
with $\mu$ gives us
\begin{align*}
\bigl[\mathsf{A}_1^{\mathsf{F}}\mathsf{T}_0 \mathsf{T}_0\bigr] \cdot \mu  =
\bigl[\mathsf{A}_1^{\mathsf{F}}\mathsf{T}_1\bigr]\cdot \mu + 2 \bigl[\mathsf{A}_1^{\mathsf{L}}\bigr]\cdot \mu.
\end{align*}
In contrast to the earlier case, the intersection of
$\bigl[\mathsf{A}_1^{\mathsf{L}}\bigr]$ with $\mu$ is nonzero (or at least not necessarily zero)
because we are making the curve pass through $\delta_d-2$ points; that is precisely the right number
of points to enumerate $1$-nodal curves, with the node lying on a line.

\begin{proof}[Proof of Theorem \ref{Tk1Tk2_etc_alt_way}]
We first prove the theorem for $n = 1$ and $k_1 = 1$, namely
we prove equation \eqref{T0T0_collides_T1_aagg}.
The main set-theoretic statement that we need to prove is as follows:
consider the component of the closure $\overline{\mathsf{T}_0 \mathsf{T}}_0$.
where the two marked points are equal. Then the first derivative of the
polynomial (defining the curve) along the direction of the line
(evaluated at the marked point) is zero. In other words, if
$(H_1, H_d, p, p) \in \overline{\mathsf{T}_0 \mathsf{T}}_0$, then
$(H_1, H_d, p) \in \overline{\mathsf{T}}_1$. Let us prove this assertion.

We continue with the set-up of the proof of Proposition \ref{prp_smth_mfld}
and Theorem \ref{theorem_for_many_Tks}.
Choose coordinate where the designated line stays the $x$-axis.
Let $f_t$ be a curve that passes through the origin and also through the point $(t, 0)$.
The expression for $f_t$ is of the form
$
f_t(x, y)  =  \varphi_t(x)+ y \mathcal{R}_t(x, y)$,
where \smash{$\varphi_{t}(x) = \bigl(f_{t_{10}} x + \frac{f_{t_{20}}}{2} x^2 + \cdots\bigr)$}.
Since the curve passes through $(t, 0)$, it follows that
$
\varphi_{t}(x)  =  (\mathcal{K}_t(x))x(x-t)
$
for some function $\mathcal{K}_t(x)$. Denote $f_0$ by $f$. Hence,
$
f(x, y) =  (\mathcal{K}_0(x))x^2 + y \mathcal{R}_0(x, y)$.
It is a simple check to see that $f_{00}$ and $f_{10}$ are both zero.
Hence, if two $\mathsf{T}_0$ points collide, then we get a point which is at least as degenerate as a
$\mathsf{T}_1$ point (it could be even more degenerate). This proves the
assertion we made.

We now prove a more general statement.
We claim the following:
consider the component of the closure $\overline{\mathsf{T}_{k} \mathsf{T}}_0$
where the two marked points are equal. Then the $(k+1)$-th derivative of the
polynomial (defining the curve) along the direction of the line
(evaluated at the marked point) is zero. In other words, if
$(H_1, H_d, p, p) \in \overline{\mathsf{T}_{k} \mathsf{T}}_0$, then
$(H_1, H_d, p) \in \overline{\mathsf{T}}_{k+1}$.

In order to prove the above assertion, put the
$\mathsf{T}_k$ point at $(0, 0)$ and the $\mathsf{T}_0$ point at $(t, 0)$.
The expression for $f_t$ is going to be of the form
\begin{align*}
f_t(x, y)& = \varphi_t(x)+ y \mathcal{R}_t(x, y),\qquad
\textnormal{where} \quad \varphi_{t}(x) =  f_{t_{10}} x + \frac{f_{t_{20}}}{2} x^2 + \frac{f_{t_{30}}}{6} x^3 + \cdots .
\end{align*}
Since the curve is tangent of the $x$-axis to order $k$ and it also passes through $(t, 0)$,
it follows that
$
\varphi_{t}(x)  =  (\mathcal{K}_t(x))x^{k+1}(x-t)
$
for some function $\mathcal{K}_t(x)$. To see what happens in the limit as~$t$ goes to zero, denote $f_0$ by $f$.
Hence,
\[
f(x, y) =  (\mathcal{K}_0(x))x^{k+2} + y \mathcal{R}_0(x, y).
\]
It is a simple check that $f_{00}, f_{10}, \dots, f_{k+1, 0}$ are all zero.
Hence, if a $\mathsf{T}_k$ and a $\mathsf{T}_0$ point collide, then
we get a point which is at least as degenerate as a $\mathsf{T}_{k+1}$ point (it could be even more degenerate).

We now need to prove the converse of the above assertion.
We claim the following: if $(H_1, H_d, p) \in \overline{\mathsf{T}}_{k+1}$,
then $(H_1, H_d, p, p) \in \overline{\mathsf{T}_{k} \mathsf{T}}_0$.

First we note that to prove the above claim, it is sufficient to prove the following:
if $(H_1, H_d, p) \in \mathsf{T}_{k+1}$,
then $(H_1, H_d, p, p) \in \overline{\mathsf{T}_{k} \mathsf{T}}_0$.
This is because if $A$ is a subset of $B$, then closure of $A$ is a subset of
closure of $B$. Hence, if $B$ is a closed set, then to show that closure of~$A$
is a subset of $B$, it is sufficient to show that $A$ is a subset of $B$.

We start by proving the claim for $k = 0$, namely, we show that
every $\mathsf{T}_1$ point can be obtained as a limit to two $\mathsf{T}_0$ points.

Let $f \in (\mathsf{T}_1)_{\textnormal{Aff}}$. This means that $f_{00}$ and $f_{10}$ are both equal to zero.
Hence, $f$ is given by
\[
f(x, y)  =  \left(\frac{f_{20}}{2} x^2 + \frac{f_{30}}{6} x^3 + \cdots\right) + y \mathcal{R}(x, y).
\]
It will be shown that there exists a point $(f_t, (t, 0))$ close to $(f, (0, 0))$, such that
$f_t$ passes through the origin and $(t, 0)$. Note that
since $f_t$ passes through the origin, it is of the form
\[
f_t(x, y) = f_{t_{10}}x+ \left(\frac{f_{t_{20}}}{2} x^2 + \frac{f_{t_{30}}}{6} x^3 +\cdots\right) +
y \mathcal{R}_t(x, y).
\]
Furthermore, $f_{t_{10}}$ has to be small (since $f_t$ is close to $f$ and $f_{10}$ is zero).
Impose the condition that $f_{t}(t, 0) = 0$. Plugging in $(t, 0)$ inside $f_t$ and using the fact
that $t \neq 0$, it follows that
\begin{align}
f_{t_{10}}&= -\frac{f_{t_{20}}}{2} t + O\bigl(t^2\bigr). \label{ft_10}
\end{align}
Hence, we have constructed this nearby curve $f_t$ and a marked point $(t, 0)$ different from the origin
that lies on the curve and the line.

To find the multiplicity of the intersection, we note that using equation \eqref{ft_10},
using the fact that $f_{t_{20}} \neq 0$
and the implicit function theorem, we can rewrite it as
\begin{align}
t & = -\frac{2}{f_{t_{20}}} f_{t_{10}} + O\bigl(f_{t_{10}}^2\bigr). \label{ft_10_t}
\end{align}
Setting the two points to be equal is the same as setting $t$ to be equal to zero.
By equation~\eqref{ft_10_t}, the order of vanishing of $t$ is one. This
justifies the multiplicity.

We now prove the assertion for a general $k$,
i.e., we show that
every $\mathsf{T}_{k+1}$ curve is in the limit of a $\mathsf{T}_{k}$
and $\mathsf{T}_0$ point.
Let $f \in (\mathsf{T}_{k+1})_{\textnormal{Aff}}$.
This means that $f_{00}, f_{10}, \dots, f_{k+1,0}$ are all equal to zero.
Hence, $f$ is given by
\[
f(x, y)  =  \left(\frac{f_{k+2,0}}{(k+2)!} x^{k+2} +
\frac{f_{k+3,0}}{(k+3)!} x^{k+3} + \cdots\right) + y \mathcal{R}(x, y).
\]
It will be shown that there exists a point $(f_t, (t, 0))$ close to $(f, (0, 0))$, such that
$f_t \in (\mathsf{T}_{k})_{\textnormal{Aff}}$
and~$(t, 0)$.
Since $f_t \in (\mathsf{T}_{k})_{\textnormal{Aff}}$,
it is of the form
\[
f_t(x, y) = \frac{f_{t_{k+1, 0}}}{(k+1)!}x^{k+1}+
\left(\frac{f_{t_{k+2, 0}}}{(k+2)!} x^{k+2} + \frac{f_{t_{k+3,0}}}{(k+3)!} x^{k+3} +\cdots\right) +
y \mathcal{R}_t(x, y).
\]
Furthermore, $f_{t_{k+1, 0}}$ has to be small (since $f_t$ is close to $f$ and $f_{k+1, 0}$ is zero).
Impose the condition that $f_{t}(t, 0) = 0$. Plugging in $(t, 0)$ inside $f_t$ and using the fact
that $t \neq 0$, it follows that
\begin{align}
f_{t_{k+1, 0}}&= -\frac{f_{t_{k+2, 0}}}{(k+2)} t + O\bigl(t^2\bigr). \label{ft_k0}
\end{align}
Hence, we have constructed this nearby curve $f_t$ and a marked point $(t, 0)$ different from the origin
that lies on the curve and the line.

To find the multiplicity of the intersection, we note that using equation \eqref{ft_k0},
using the fact that $f_{t_{k+2, 0}} \neq 0$ (which is true because after we make the curves
pass through generic points)
and the implicit function theorem, we can rewrite it as
\begin{align}
t & = -\frac{2}{f_{t_{k+2, 0}}} f_{t_{k+1, 0}} + O\bigl(f_{t_{k+1, 0}}^2\bigr). \label{ft_k0_t}
\end{align}
Setting the two points to be equal is the same as setting $t$ to be equal to zero.
By equation~\eqref{ft_k0_t}, the order of vanishing of $t$ is one. This
justifies the multiplicity.

We summarize the set-theoretic statement we have just proved.
We have shown that $(H_1, H_d,\allowbreak  p, p)$ belongs to $\overline{\mathsf{T}_k \mathsf{T}}_0$
if and only if $(H_1, H_d, p)$ belongs to $\overline{\mathsf{T}}_{k+1}$.

We now examine what happens when there are more than two points involved.
We explain with the help of a single example; the general case follows in a similar way.
Consider the following assertion:
\[
[\mathsf{T}_2 \mathsf{T}_0 \mathsf{T}_0] \cdot \bigl[\Delta^{\mathsf{L}}_{23}\bigr]  =
\pi^*[\mathsf{T}_2\mathsf{T}_1]\cdot [\Delta_{23}].
\]
The closure claim that we need to prove is as follows:
$(H_1, H_d, q, p, p)$ belongs to \smash{$\overline{\mathsf{T}_2\mathsf{T}_0 \mathsf{T}}_0$}
if and only if $(H_1, H_d, p)$ belongs to~\smash{$\overline{\mathsf{T}}_{1}$}.
One direction is the same as before namely that if~${(H_1, H_d, q, p, p)}$ belongs to
$\overline{\mathsf{T}_2\mathsf{T}_0 \mathsf{T}}_0$, then~$(H_1, H_d, p)$ belongs to $\overline{\mathsf{T}}_{1}$. The fact that
$(H_1, H_d, q)$ belongs to $\overline{\mathsf{T}}_2$ makes no difference in the proof.
It is the converse that requires a little bit more argument. We need to show that
every $\mathsf{T}_2 \mathsf{T}_1$ point can be obtained as a limit of
$\mathsf{T}_2 \mathsf{T}_0 \mathsf{T}_0$. First see what we did when we
had to show every $\mathsf{T}_1$ point can be obtained as a limit of
$\mathsf{T}_0 \mathsf{T}_0$. We constructed the curve as given by equation
\eqref{ft_10}. The problem with equation \eqref{ft_10} is that if we set
$f_{t_{10}}$ as given by \eqref{ft_10} and the remaining $f_{t_{ij}}$ to be
complex numbers close to $f_{ij}$, then this curve does not satisfy the
$\mathsf{T}_2$ condition at $q$. The problem is the \smash{$f_{t_{ij}}$} are not all free.
However, we have shown through the proof of Proposition~\ref{prp_smth_mfld} that
restricted to the submanifold~$\mathsf{T}_2 \overline{\mathsf{T}}_0$, the section induced by $f_{10}$ is transverse to
zero. Hence, $f_{10}$ is an actual coordinate on the submanifold~$\mathsf{T}_2 \overline{\mathsf{T}}_0$.
Hence, \eqref{ft_10} defines a curve lying in~$\mathsf{T}_2 \overline{\mathsf{T}}_0$ that
converges to an element of~$\mathsf{T}_2 \overline{\mathsf{T}}_1$ as~$t$ goes to zero.
The multiplicity computation is the same. The general case follows from equation~\eqref{ft_k0} and using the fact that $f_{k+1, 0}$ is a~coordinate on~$\mathsf{T}_2 \overline{\mathsf{T}}_{k}$. When there are more than three points involved, the
same argument holds via the proof of Proposition~\ref{prp_smth_mfld},
namely that~$f_{k_n, 0}$ is a coordinate on
$\overline{\mathsf{T}}_{k_1} \overline{\mathsf{T}}_{k_2}
\dots \overline{\mathsf{T}}_{k_{n-1}}\overline{\mathsf{T}}_{k_{n}-1}$.
This completes the proof of Theorem~\ref{Tk1Tk2_etc_alt_way}.
 \end{proof}
\begin{thebibliography}{99}
\bibitem[11]{GH3}
Griffiths P., Harris J., Principles of algebraic geometry, \textit{Wiley Classics Lib.},
  \href{https://doi.org/10.1002/9781118032527}{John Wiley \& Sons}, Inc., New York, 1994.

\end{thebibliography}
\end{document}
